\infographicchapter{Magnetoencephalography (MEG)}{meg.png}
\label{ch:infographic-meg}

\section{Overview}

Magnetoencephalography (MEG) is a completely painless, silent, noninvasive brain test that listens to the tiny natural magnetic fields created by active brain cells. Whenever thousands of nerve cells in the outer brain fire together, they produce faint electrical currents and matching magnetic fields---millions of times weaker than the Earth's magnetism. Unlike electrical signals measured on the scalp by EEG, magnetic fields pass smoothly through the skull and scalp without blurring. When doctors lay MEG brain-activity maps onto a sharp structural MRI scan (a combination called magnetic source imaging, or MSI), they can pinpoint brain activity to within a few millimeters and track its timing down to the millisecond.

Traditional whole-head MEG machines use hundreds of ultra-sensitive superconducting sensors (called SQUIDs) cooled with liquid helium inside a smooth helmet, all housed inside a thick magnetically shielded room that blocks outside noise from elevators, traffic, and power lines. A newer breakthrough uses lightweight, wearable sensors called optically pumped magnetometers (OPMs), which work at room temperature and fit snugly like a cap---making MEG much easier for babies, young children, and people who have trouble sitting still.

In clinical care, MEG is a well-established tool for patients with hard-to-treat focal epilepsy who are being evaluated for surgery. It pinpoints where abnormal seizure sparks start, clarifies unclear brain workups, and guides surgeons on exactly where to place monitoring electrodes. Doctors also use MEG before brain tumor or epilepsy surgery to map vital areas for touch, movement, hearing, vision, and speech so surgeons can remove disease safely while protecting everyday abilities.

\section{Underlying Physical Principles}
\begin{itemize}
  \item \textbf{Natural Magnetism From Synaptic Currents:} When roughly 50,000 neighboring cortical pyramidal cells receive messages together, the slow electrical currents flowing along their branches (postsynaptic potentials) act like miniature batteries---whereas brief firing spikes largely cancel out---producing a weak magnetic field right outside the head.
  \item \textbf{Passing Cleanly Through the Skull:} Because the skull, scalp, and brain fluid do not distort magnetic fields the way they smear electrical signals, MEG captures much cleaner, less blurred brain signals than standard scalp EEG---even in people who have had prior skull surgery.
  \item \textbf{Best at Reading Brain Folds:} Because of the round shape of the head, MEG is especially sensitive to brain currents flowing sideways along the walls of the brain's natural folds (sulci), while currents pointing straight up toward the scalp produce weaker external magnetic signals.
  \item \textbf{Ultra-Sensitive Quiet Detectors:} Inside a heavily shielded room, either liquid-helium-cooled SQUID helmets or room-temperature wearable OPM caps detect these ultra-faint magnetic waves without sending any energy into your head.
  \item \textbf{Mapping Signals Onto Your Brain MRI:} Specialized computer programs calculate where inside the brain the magnetic signals came from and line them up precisely with your own 3D brain MRI using small forehead and ear position markers.
\end{itemize}

\section{Imaging Mechanism}
\begin{enumerate}
  \item \textbf{Preparing Position Markers and Entering the Quiet Room:} Outside the shielded room, a technologist places small sticky head-position coils on your forehead and behind your ears, adds heart and eye-blink reference stickers, and gently traces a few face landmarks with a wand so the MEG will match your MRI. Inside the shielded room, you sit or recline with your head resting inside the smooth helmet (or wearing an OPM cap).
  \item \textbf{Silent Listening and Simple Tasks:} While the position markers gently track any tiny head movements, hundreds of sensors silently listen to your resting brainwaves or record brain responses while you feel light finger taps, hear soft tones, view pictures, or think of words.
  \item \textbf{Amplifying and Digitizing the Brain Signals:} Sensitive electronics boost the ultra-faint magnetic signals, filter out stray electrical noise, and record them thousands of times per second.
  \item \textbf{Cleaning Artifacts and Pinpointing Sources:} Computer software filters out heartbeat, eye-blink, and outside magnetic interference, then calculates the exact 3D brain spots creating the signals and aligns them onto your structural MRI.
  \item \textbf{Creating Surgical Brain Maps:} Clusters of seizure sparks and vital speech, movement, or sensory zones appear as color-coded 3D overlays on your MRI slices, ready for your epilepsy team and surgical navigation screens.
\end{enumerate}

\section{Key Technical Parameters}
\noindent
\begin{tabularx}{\textwidth}{@{}l X@{}}
\toprule
\textbf{Parameter} & \textbf{Typical Value / Description} \\
\midrule
Physical Principle & Passive listening to ultra-faint natural magnetic fields created by firing nerve cells in the outer brain \\
Energy / Signal Source & Natural brain-cell currents (no energy sent into the body); recorded by cooled SQUID helmets or wearable OPM sensors \\
Typical Frequency / Energy & Ultra-faint 10--1,000 femtotesla magnetic signals across 0.1--100 Hz brainwave bands \\
Spatial Resolution & 2--10 mm accuracy on the brain surface when matched to a 3D MRI; millisecond timing speed \\
Typical Scan Duration & 60--90 minutes in short 5--15 minute blocks (plus 20--30 minutes for sensor and marker setup) \\
Contrast Agent & None (completely passive and dye-free) \\
Ionizing Radiation & None \\
\bottomrule
\end{tabularx}

\section{Common Symptoms That Prompt This Scan}
\begin{itemize}
  \item \textbf{Seizures That Do Not Respond to Medication:} Frequent focal seizures---such as sudden staring spells, unusual smells, jerking of one limb, or brief confusion---when doctors are evaluating epilepsy surgery.
  \item \textbf{Seizures Returning After Prior Brain Surgery:} Recurrent seizure episodes after an earlier operation, where skull openings distort standard scalp EEG but do not blur MEG signals.
  \item \textbf{Brain Tumor Symptoms Near Speech or Movement Centers:} Headaches, subtle speech trouble, or limb tingling from a brain tumor or vessel tangle when surgeons need precise noninvasive functional mapping.
  \item \textbf{Note---Specialized Referral Test:} MEG is not a routine first test for a single seizure or headache; it is ordered by epilepsy and neurosurgery teams at specialized centers.
\end{itemize}

\section{Clinical Usages}

\subsection{Primary Usages}
\begin{itemize}
  \item \textbf{Finding Where Hard-to-Treat Seizures Start:} Pinpointing abnormal seizure sparks between seizures in people with medication-resistant focal epilepsy---especially when standard MRI looks normal or other tests disagree.
  \item \textbf{Guiding Surgical Monitoring Electrodes:} Narrowing down the suspected seizure zone with high accuracy so neurosurgeons know exactly where to place internal monitoring wires (stereo-EEG or grids), improving chances of seizure freedom.
  \item \textbf{Mapping Vital Brain Functions Before Surgery:} Painlessly locating touch, movement, hearing, and vision centers before removing brain tumors, blood-vessel tangles, or seizure areas so critical abilities are protected.
  \item \textbf{Mapping Speech and Language:} Helping determine which side of the brain controls language and where speaking and understanding areas sit, offering a gentle noninvasive partner or alternative to invasive Wada testing and fMRI.
\end{itemize}

\subsection{Secondary Usages}
\begin{itemize}
  \item \textbf{Checking Returning Seizures After Prior Brain Surgery:} Because surgical bone openings distort scalp EEG much more than magnetic fields, MEG is especially helpful when seizures come back after a past brain operation.
  \item \textbf{Young Children and Movement-Prone Patients:} Emerging wearable OPM caps avoid rigid helmets and cold tanks, making MEG much more feasible for infants, young children, and people with autism or movement disorders.
  \item \textbf{Concussion and Memory Research:} Investigational studies use MEG to look for abnormal slow brainwaves and disrupted brain networks in mild traumatic brain injury, PTSD, Alzheimer's disease, and Parkinson's disease.
  \item \textbf{Brain Science and Brain--Computer Interfaces:} Research centers use MEG's millisecond speed to study how brain networks talk to each other in real time.
\end{itemize}

\section{Advantages}
\begin{itemize}
  \item \textbf{Direct Millisecond Timing:} Measures active brain-cell firing directly in real time---not delayed blood flow---tracking rapid seizure spread and fast thinking processes.
  \item \textbf{Not Blurred by the Skull:} Magnetic fields pass cleanly through the skull, scalp, and prior surgical bone openings, avoiding the smearing that limits scalp EEG.
  \item \textbf{Sharp Surface Brain Mapping:} Paired with a 3D brain MRI, MEG pinpoints seizure zones and vital function areas to within millimeters.
  \item \textbf{Completely Painless and Radiation-Free:} Uses zero X-rays, zero contrast dye, and zero applied magnetic fields, making it safe for children and repeated tests.
  \item \textbf{Growing Comfort With Wearable Sensors:} Newer room-temperature OPM caps remove the need for rigid helmets, improving comfort for children and restless patients.
\end{itemize}

\section{Disadvantages}
\begin{itemize}
  \item \textbf{Less Sensitive to Deep Brain Areas:} Magnetic signals fade quickly with distance, so MEG sees the outer brain surface much better than deep inner structures or currents pointing straight toward the scalp.
  \item \textbf{Very Sensitive to Metal and Outside Noise:} Because brain magnetism is so faint, braces, magnetized dental work, vagus nerve stimulators, metal implants, or head movement can easily drown out the signal.
  \item \textbf{Needs a Separate Brain MRI:} MEG shows brain activity, not anatomy, so you also need a high-detail structural MRI to see exactly where the signals sit. Mathematical modeling also requires expert interpretation.
  \item \textbf{High Cost and Rare Availability:} Traditional SQUID systems require costly shielded rooms, liquid helium, and specialized teams, so MEG is available mainly at major academic epilepsy centers.
  \item \textbf{Short Recording Window:} A 1--2 hour visit mainly captures quiet between-seizure sparks and can miss rare seizure activity compared with multi-day hospital video-EEG stays.
\end{itemize}

\section{Contraindications}

\subsection{Absolute}
\begin{itemize}
  \item \textbf{None for MEG Itself:} Passive MEG has no strict medical exclusions, because the machine sends no magnetic fields, radio waves, or radiation into your body.
\end{itemize}

\subsection{Relative}
\begin{itemize}
  \item \textbf{Heavy Metal Dental Work or Active Implants:} Magnetized orthodontic braces, extensive metal dental work, cranial hardware, shunts, or active vagus nerve stimulators can create strong magnetic noise that makes MEG difficult or impossible to read clearly.
  \item \textbf{Claustrophobia or Uncontrolled Movement:} Severe anxiety inside the closed shielded room, strong tremors, or frequent convulsive seizures may limit image quality unless calming support or wearable OPM sensors are available.
  \item \textbf{Inability to Get the Matching MRI:} Because clinical MEG maps are layered onto a brain MRI, patients with MRI-unsafe implants cannot complete the MRI half and must use CT instead with less anatomical detail.
\end{itemize}

\section{Patient Safety Guidelines}
\begin{itemize}
  \item \textbf{Removing All External Metal:} Before entering the shielded room, you remove all metallic clothing, jewelry, hairpins, and accessories, and let staff know about braces, shunts, or vagus nerve stimulators so they can plan noise filtering.
  \item \textbf{Seizure Safety Support:} Because many patients are being checked for epilepsy---sometimes after staying up late---trained epilepsy staff watch you continuously on video and audio with emergency airway and rescue medicines right at hand.
  \item \textbf{Easy Room Exit and Comfort:} The heavy shielded-room door opens quickly from both inside and outside, a two-way intercom and call button keep you in constant touch, and a parent or caregiver can often accompany anxious or young patients.
  \item \textbf{Equipment Safety Checks:} Cooled SQUID suites have oxygen monitors and safety vents for liquid helium; wearable OPM caps are monitored for gentle surface warmth and comfortable fit.
  \item \textbf{MRI Safety Screening for the Companion Scan:} Full standard MRI metal screening is completed before the separate structural brain MRI needed to build your final map.
\end{itemize}

\section{Prerequisite Tests}
\begin{itemize}
  \item \textbf{No Blood Tests or Routine Fasting:} MEG is completely passive and dye-free, so no lab work is required, and most adults need no fasting---though your care team will tell you if food or drink must be limited ahead of time (especially if sedation is planned).
  \item \textbf{High-Detail Brain MRI:} A recent sharp 3D brain MRI is required before or alongside MEG so your brain activity can be mapped onto your exact brain folds.
  \item \textbf{Prior Epilepsy and Neurology Workup:} Before pre-surgery MEG, your epilepsy team reviews your seizure history, neurological exam, and scalp EEG (or hospital video-EEG) to focus the test.
  \item \textbf{Metal and Function Check:} Staff document any dental hardware or implanted devices beforehand, and do quick baseline checks of speech, movement, hearing, or vision if functional mapping is planned.
\end{itemize}

\section{Preparing for the Scan}
\begin{itemize}
  \item \textbf{Clean Hair and Bare Skin:} Wash your hair the night before or morning of the test with plain shampoo. Skip hair gel, hairspray, oils, conditioners, makeup, mascara, or metallic skincare creams, because tiny metal particles in products can distort the ultra-sensitive sensors.
  \item \textbf{Metal-Free Clothing:} Change into a hospital gown or metal-free scrubs and remove all jewelry, watches, glasses, hairpins, belts, underwire bras, and removable dental work. Tell staff about permanent braces, fillings, or a vagus nerve stimulator.
  \item \textbf{Sleep and Seizure-Medicine Rules:} Keep taking all prescribed seizure and regular medicines unless your epilepsy doctor explicitly orders a supervised change. If your team asks you to sleep less the night before (to encourage subtle seizure sparks), follow their exact instructions, avoid caffeine, and arrange a safe ride.
  \item \textbf{Position Coils and Landmark Tracing:} Before entering the shielded room, a technologist tapes small position coils to your forehead and behind your ears, places heart and eye-blink stickers (plus scalp EEG leads if ordered), and gently traces landmarks on your nose and ears.
  \item \textbf{Practicing Tasks and Getting Comfortable:} You will practice any simple finger-tapping, listening, or picture-naming tasks beforehand, use the restroom before the 1--2 hour session, and see how the intercom and door work so you feel completely at ease.
  \item \textbf{Anxiety Support, Pediatric Sedation, and Visit Timing:} If having equipment placed closely around your head makes you anxious, ask your care team beforehand about a mild sedative; infants and young children often receive sedation or general anesthesia to help them stay completely still during traditional helmet scans. Plan for up to 3--4 hours at the clinic overall (including setup and breaks), even though the actual scanning rarely lasts more than two hours.
\end{itemize}

\section{During Scan}
\begin{itemize}
  \item \textbf{Quiet Helmet Positioning:} Inside the quiet shielded room, you sit in a comfortable chair or recline on a bed with your head resting gently inside the smooth helmet (or wearing a light OPM cap). Keeping your head, jaw, and shoulders relaxed and still keeps the readings clean.
  \item \textbf{Resting and Simple Tasks:} Breathe naturally and try not to blink heavily, clench your jaw, or swallow repeatedly during active blocks. During resting epilepsy checks, you may be encouraged to relax or lightly doze; during function mapping, you may feel gentle finger taps, hear tones, view pictures, or silently name objects.
  \item \textbf{Completely Silent and Painless:} Unlike MRI, MEG makes no loud knocking noises and uses no needles or injections. You will not feel the magnetic sensors at all---only the light tape on your forehead or gentle task cues.
  \item \textbf{Staying in Touch With Staff:} The technologist watches you continuously on video and talks with you through an intercom, and the door can be opened right away at any time. Speak up immediately if you feel a seizure warning (aura), dizziness, neck stiffness, or anxiety.
  \item \textbf{Comfortable Pacing:} Recording is broken into easy 5--15 minute blocks with short rest breaks to stretch or adjust, taking about 60--90 minutes inside the quiet room.
\end{itemize}

\section{Normal Appearance}
\begin{itemize}
  \item \textbf{Balanced Natural Brain Rhythms:} A normal awake MEG shows smooth, even background rhythms on both sides of the brain with no localized slow-wave patches.
  \item \textbf{No Clustered Seizure Sparks:} A normal epilepsy recording shows no repeated, tightly grouped clusters of abnormal magnetic spikes.
  \item \textbf{Expected Function Locations:} During brain mapping, touch, movement, hearing, vision, and speech responses light up cleanly in their normal anatomical neighborhoods.
  \item \textbf{Crisp Alignment With MRI:} The forehead and ear position markers line up smoothly with the scalp outline on your 3D brain MRI.
\end{itemize}

\section{Observing for Abnormal Patterns}
\begin{itemize}
  \item \textbf{Tightly Clustered Seizure Sparks:} Repeated abnormal spikes that cluster tightly in one spot on your co-registered brain MRI pinpoint the irritated brain region where focal seizures are likely starting.
  \item \textbf{Localized Slow-Wave Zones:} Patches of unusually slow brainwaves often surround brain tumors, old strokes, developmental brain-fold quirks, or areas bruised by head injury.
  \item \textbf{Shifted Vital Brain Areas:} If touch, movement, hearing, or speech responses appear shifted away from their usual landmarks, it alerts surgeons that a tumor or birth condition has pushed or reorganized vital brain tissue.
  \item \textbf{Delayed or Weakened Responses:} Slower or weaker magnetic responses during sensory or speech tasks can point to pathway injury, tumor involvement, or which side of the brain leads language.
  \item \textbf{Harmless Body and Metal Look-Alikes:} Heartbeats, eye blinks, jaw clenching, magnetized dental work, or vagus nerve stimulators create strong magnetic blips that can mimic seizure spikes unless filtered out with special software and heart/eye stickers.
\end{itemize}

\section{Dealing with Uncertainties}
\begin{itemize}
  \item \textbf{Deep Brain and Straight-Up Signals:} Because magnetic signals fade with distance and MEG sees sideways brain-fold currents best, seizure spots deep in the center of the brain can produce faint or missing MEG readings.
  \item \textbf{Quiet Recording Days:} Because a MEG session lasts only 60--90 minutes, rare seizure sparks may simply not happen during the visit. A ``normal'' MEG never rules out epilepsy.
  \item \textbf{Body and Metal Noise:} Dental hardware, vagus nerve stimulators, heartbeats, eye blinks, jaw clenching, and head shifting can blur or mimic real brain spikes.
  \item \textbf{Modeling and Alignment Limits (The Inverse Problem):} Pinpointing 3D brain spots from outside sensors has no single unique mathematical answer, so different models or tiny alignment shifts can nudge estimated locations by a few millimeters. Colored activity overlays on your MRI therefore show statistical probability maps rather than exact anatomical borders.
  \item \textbf{Team Review With Multiple Tests:} Specialists clean up noise with advanced filtering software, sometimes repeat a sleep-deprived or OPM scan, and always combine MEG with video-EEG, 3 Tesla MRI, PET scans, and surgical depth electrodes when needed.
\end{itemize}

\section{Augmenting Methods}
\begin{itemize}
  \item \textbf{High-Detail 3 Tesla Brain MRI:} Provides the sharp anatomical backdrop onto which MEG spots are layered, helping uncover subtle brain-fold abnormalities or hippocampal scarring.
  \item \textbf{Simultaneous Scalp EEG and Hospital Video-EEG:} Captures the straight-up brain signals MEG sees less easily and records actual seizures over several days to match brainwaves with physical symptoms.
  \item \textbf{FDG PET and Seizure SPECT Scans:} Show areas of low brain energy use between seizures (PET) or surges of blood flow during a seizure (SPECT), offering independent confirmation when MRI looks normal.
  \item \textbf{Functional MRI, Fiber Tracking, and Memory Testing:} Complement MEG language and movement maps by showing deep brain wiring (DTI) and detailed thinking and memory strengths.
  \item \textbf{Surgical Depth Electrodes (Stereo-EEG) and Awake Mapping:} In complex surgery cases, MEG clusters guide exactly where neurosurgeons place depth electrodes or test brain areas directly during surgery.
\end{itemize}

\section{Comparison with Alternative Modalities}
\begin{itemize}
  \item \textbf{Compared With Scalp EEG:} MEG is barely blurred by the skull (especially helpful after prior skull surgery) and pinpoints surface spots more tightly, while scalp EEG is far less costly, portable, and practical for multi-day hospital seizure recording.
  \item \textbf{Compared With Functional MRI (fMRI):} MEG reads brain-cell firing directly at millisecond speed rather than slower blood-flow changes, whereas fMRI sees deep inside the entire brain and is much more widely available.
  \item \textbf{Compared With PET and SPECT:} MEG pinpoints brief seizure sparks sharply with zero radiation, whereas PET and SPECT map broader energy or blood-flow zones and reach deeper brain structures more easily.
  \item \textbf{Compared With Invasive Stereo-EEG:} MEG is completely painless and surveys the whole outer brain at once to form a surgical plan, whereas stereo-EEG directly samples deep brain areas but requires surgery to place thin wires.
\end{itemize}

\section{Follow-up Steps}
\begin{itemize}
  \item \textbf{Immediate Return to Your Day:} Because MEG is completely passive and noninvasive, most people resume normal eating and regular activities right away. If you or your child received a sedative, stayed up late (sleep deprivation), or adjusted seizure medicines, plan to wait a couple of hours for sedation to wear off, arrange a safe ride home, and avoid driving until fully rested and cleared by your doctor.
  \item \textbf{MRI Mapping and Epilepsy Team Review:} Your MEG seizure clusters and function maps are layered onto your 3D brain MRI and reviewed together by epilepsy doctors, neuroradiologists, neurophysiologists, and neurosurgeons alongside your EEG, PET, and memory tests.
  \item \textbf{Guiding Safe Brain Surgery:} Clear MEG results help surgeons plan targeted depth-electrode placement or guide tumor and epilepsy surgery safely around vital speech and movement zones.
  \item \textbf{Long-Term Tracking:} After surgery, your team tracks your seizure control and recovery; if seizures ever return after skull surgery, MEG can be repeated cleanly without skull-opening distortion.
\end{itemize}

\section{Key References}
\begin{itemize}
  \item Bagic AI, Knowlton RC, Rose DF, Ebersole JS; ACMEGS Clinical Practice Guideline Committee. American Clinical Magnetoencephalography Society Clinical Practice Guideline 1: Recording and analysis of spontaneous cerebral activity. \textit{Journal of Clinical Neurophysiology}. 2011;28(4):348--354.
  \item Hari R, Baillet S, Barnes G, et al. IFCN-endorsed practical guidelines for clinical magnetoencephalography (MEG). \textit{Clinical Neurophysiology}. 2018;129(8):1720--1747.
  \item Burgess RC, Funke ME, Bowyer SM, et al. American Clinical Magnetoencephalography Society Clinical Practice Guideline 2: Presurgical functional brain mapping using magnetic evoked fields. \textit{Journal of Clinical Neurophysiology}. 2011;28(4):355--361.
  \item Brookes MJ, Leggett J, Rea M, et al. Magnetoencephalography with optically pumped magnetometers (OPM-MEG): the next generation of functional neuroimaging. \textit{Trends in Neurosciences}. 2022;45(8):621--634.
  \item Mayo Clinic. \textit{Magnetoencephalography (MEG)}. Mayo Clinic. Available at: \url{https://www.mayoclinic.org/tests-procedures/magnetoencephalography-meg/about/pac-20575365}.
\end{itemize}

\section{Further Reading}
\begin{itemize}
  \item Wikipedia contributors. \textit{Magnetoencephalography}. Wikipedia, The Free Encyclopedia. Available at: \url{https://en.wikipedia.org/wiki/Magnetoencephalography}.
\end{itemize}
